\documentclass[11pt,a4paper]{article}
\usepackage[spanish]{babel}
\usepackage[T1]{fontenc}

\usepackage[margin=2cm,headheight=14pt,headsep=18pt,footskip=22pt]{geometry}
\usepackage[table]{xcolor}
\usepackage{tabularx}
\usepackage{longtable}
\usepackage{fancyhdr}
\usepackage{lastpage}
\usepackage{titlesec}
\usepackage{lmodern}

\renewcommand{\familydefault}{\sfdefault}
\definecolor{VerdeUteq}{HTML}{006537}
\definecolor{VerdeTabla}{HTML}{11753C}
\definecolor{AzulTitulo}{HTML}{19598A}
\definecolor{Dorado}{HTML}{D89D0B}

\newcommand{\estudiante}{Jonathan Holguín}
\newcommand{\asignatura}{Ingeniería de Requerimientos (ISR-401)}

\pagestyle{fancy}
\fancyhf{}
\fancyhead[L]{\small UTEQ \textperiodcentered\ \asignatura}
\fancyhead[R]{\small FORM-TL-01 \textperiodcentered\ SPA 2026--2027}
\fancyfoot[L]{\small \estudiante}
\fancyfoot[R]{\small Página \thepage\ de \pageref{LastPage}}
\renewcommand{\footrulewidth}{0.4pt}

\titleformat{\section}
  {\large\bfseries}{\thesection.}{0.9em}{\color{VerdeUteq}}

\newcommand{\cabeceraTabla}[1]{\textcolor{white}{\textbf{#1}}}

\begin{document}

\begin{center}
  {\Large\bfseries\color{VerdeUteq} Universidad Técnica Estatal de Quevedo\par}
  \vspace{0.3em}
  {\small Facultad de Ciencias de la Computación
    \textcolor{Dorado}{\textbullet} Carrera de Software\par}
  \vspace{1em}
  {\large\bfseries\color{AzulTitulo} ACTIVIDAD FORMATIVA\par}
  \vspace{0.25em}
  {\large\bfseries FORM-TL-01\par}
  {\Large\bfseries\color{VerdeUteq} Taller de tipos de requisitos\par}
\end{center}
\noindent\textcolor{Dorado}{\rule{\linewidth}{0.8pt}}

\section{Datos generales}

{\renewcommand{\arraystretch}{1.2}
\rowcolors{2}{white}{gray!12}
\noindent\begin{tabularx}{\linewidth}{|p{0.26\linewidth}|X|}
  \hline
  \rowcolor{VerdeTabla}
  \cabeceraTabla{Campo} & \cabeceraTabla{Detalle} \\ \hline
  Estudiante & \estudiante \\ \hline
  Asignatura & \asignatura \\ \hline
  Caso & Plataforma de gestión de citas médicas \\ \hline
\end{tabularx}}

\section{Clasificación de los enunciados}

{\small
\setlength{\tabcolsep}{4pt}
\renewcommand{\arraystretch}{1.1}
\begin{longtable}
  {|>{\raggedright\arraybackslash}p{0.07\linewidth}
   |>{\raggedright\arraybackslash}p{0.24\linewidth}
   |>{\raggedright\arraybackslash}p{0.155\linewidth}
   |>{\raggedright\arraybackslash}p{0.095\linewidth}
   |>{\raggedright\arraybackslash}p{0.10\linewidth}
   |>{\raggedright\arraybackslash}p{\dimexpr0.34\linewidth-12\tabcolsep-7\arrayrulewidth\relax}|}
  \hline
  \rowcolor{VerdeTabla}
  \cabeceraTabla{Código} & \cabeceraTabla{Enunciado} & \cabeceraTabla{Tipo} &
  \cabeceraTabla{Producto o proceso} & \cabeceraTabla{Sistema o software} &
  \cabeceraTabla{Justificación breve} \\ \hline
  \endhead

  E01 & El paciente podrá reservar una cita eligiendo especialidad, médico,
  fecha y hora entre los horarios disponibles.
  & Funcional & Producto & Software &
  Reservar un turno es una función de la aplicación que se entrega. El
  paciente elige un horario libre y el software guarda la cita. \\ \hline

  \rowcolor{gray!12}
  E02 & La plataforma enviará al paciente un recordatorio por mensaje de texto
  24 horas antes de su cita.
  & Funcional & Producto & Software &
  Enviar el recordatorio es una función automática de la plataforma en uso.
  Las 24 horas indican cuándo el software debe iniciar el envío. \\ \hline

  E03 & La aplicación debe ser rápida.
  & \textbf{Calidad}\newline
  Eficiencia de desempeño\newline
  Comportamiento temporal
  & Producto & Software &
  Entiendo \textbf{rápida} como una cualidad de la aplicación en uso, ligada
  a su tiempo de respuesta, aunque no se especifica qué operación se medirá.
  \\ \hline

  \rowcolor{gray!12}
  E04 & El recepcionista podrá registrar la llegada del paciente que se acerca
  a la ventanilla.
  & Funcional & Producto & Software &
  Registrar la llegada es una función de la aplicación que se entrega. El
  recepcionista la usa para dejar constancia de que el paciente llegó a
  ventanilla. \\ \hline

  E05 & La plataforma se desarrollará en Java 21 con base de datos PostgreSQL,
  tecnologías que domina el área de TI.
  & Restricción tecnológica & Producto & Software &
  Java 21 y PostgreSQL son obligatorios en la plataforma que se entrega. Esa
  condición limita la elección del lenguaje y la base de datos del software. \\ \hline

  \rowcolor{gray!12}
  E06 & El médico podrá consultar la lista de pacientes del día y marcar cada
  consulta como atendida.
  & Funcional & Producto & Software &
  Consultar y marcar consultas son funciones del producto que el médico usa
  directamente en la aplicación. \\ \hline

  E07 & Los datos de salud de los pacientes deben estar seguros.
  & \textbf{Calidad}\newline Seguridad\newline Confidencialidad
  & Producto & Software &
  Es una cualidad del producto en uso: el software protege la
  confidencialidad al limitar quién puede consultar los datos de salud. \\ \hline

  \rowcolor{gray!12}
  E08 & La aplicación móvil debe funcionar en teléfonos con Android 10 o superior.
  & Restricción técnica & Producto & Software &
  Es una restricción del producto: fija la versión mínima de Android que
  debe admitir la aplicación móvil. \\ \hline

  E09 & Una cita cancelada con menos de cuatro horas de anticipación se
  registrará como inasistencia.
  & Funcional & Producto & Software &
  Es una regla que ejecuta el software en uso: al cancelar la cita,
  comprueba cuánto falta y la registra como inasistencia si quedan menos de
  cuatro horas. \\ \hline

  \rowcolor{gray!12}
  E10 & La plataforma debe ser fácil de usar para los pacientes adultos mayores.
  & \textbf{Calidad}\newline Capacidad de interacción\newline Operabilidad
  & Producto & Software &
  Es una cualidad de la aplicación entregada: su interfaz debe ser fácil de
  manejar para los adultos mayores al hacer tareas como reservar una cita. \\ \hline
\end{longtable}}

\section{Cuantificación de los requisitos de calidad}

\section{Comparación con un compañero}

\end{document}
